% GENERATED FILE. Edit canonical lessons, then run npm run generate:books.
% canonical-source-hash: fnv1a64:1257dec488860b14
% canonical-lessons: TE-C138-vegam, TE-C138-andamaina, TE-C138-lotu, TE-C138-vedalpu, TE-C138-iruku

\chapter{Fast, Beautiful, Deep, Wide, Narrow}
\label{ch:te-fast-beautiful-deep-wide-narrow}
\hlchaptermodality{138}

\begin{chapteropening}
\textbf{By the end of this chapter:} \emph{I can say fast, beautiful, deep, wide and narrow.}

Complete the last lesson of chapter 138: I can say fast, beautiful, deep, wide and narrow.
\end{chapteropening}

\section[\texorpdfstring{\te{వేగం} (vēgaṁ)}{వేగం (vēgaṁ)}]{\te{వేగం} (vēgaṁ) — fast; speed}

\label{lesson:TE-C138-vegam}



\begin{quote}
Before the new one: say the Telugu for clean, then the Telugu for dirty.
\end{quote}

\subsection*{You'll want to know: \te{వేగం}}
\textbf{\te{వేగం}} — \emph{vēgaṁ} — \textquotedblleft{}fast; speed\textquotedblright{}.

\begin{grammarlens}[title={What you've built}]
One more word for this chapter.
\end{grammarlens}

\subsection*{Your turn}
\begin{itemize}
\raggedright
  \item \emph{Say it:} \emph{vēgaṁ}
  \item \emph{Say it:} \emph{vēgaṁ}, once more
  \item \emph{Say it:} say \emph{vēgaṁ}
  \item \emph{Recall:} say the Telugu for clean, then the Telugu for dirty, then say \emph{vēgaṁ} again
\end{itemize}

\begin{tcolorbox}[breakable,colback=teal!4,colframe=teal!35!black,title={Before you move on}]
What does \te{వేగం} mean? (\textquotedblleft{}Fast; speed\textquotedblright{}.) Say it once more.
\end{tcolorbox}
\section[\texorpdfstring{\te{అందమైన} (andamaina)}{అందమైన (andamaina)}]{\te{అందమైన} (andamaina) — beautiful}

\label{lesson:TE-C138-andamaina}



\begin{quote}
Before the new one: say the Telugu for dirty, then the Telugu for fast; speed.

An adjective goes before the noun and does not change: \textbf{\te{పాత} \te{పుస్తకం}}, \textbf{\te{కొత్త} \te{పుస్తకం}}.
\end{quote}

\subsection*{You'll want to know: \te{అందమైన}}
\textbf{\te{అందమైన}} — \emph{andamaina} — \textquotedblleft{}beautiful\textquotedblright{}.

\begin{grammarlens}[title={What you've built}]
One more word for this chapter.
\end{grammarlens}

\subsection*{Your turn}
\begin{itemize}
\raggedright
  \item \emph{Say it:} \emph{andamaina}
  \item \emph{Say it:} \emph{andamaina}, once more
  \item \emph{Say it:} say \emph{andamaina}
  \item \emph{Recall:} say the Telugu for dirty, then the Telugu for fast; speed, then say \emph{andamaina} again
\end{itemize}

\begin{tcolorbox}[breakable,colback=teal!4,colframe=teal!35!black,title={Before you move on}]
What does \te{అందమైన} mean? (\textquotedblleft{}Beautiful\textquotedblright{}.) Say it once more.
\end{tcolorbox}
\section[\texorpdfstring{\te{లోతు} (lōtu)}{లోతు (lōtu)}]{\te{లోతు} (lōtu) — deep}

\label{lesson:TE-C138-lotu}



\begin{quote}
Before the new one: say the Telugu for fast; speed, then the Telugu for beautiful.
\end{quote}

\subsection*{You'll want to know: \te{లోతు}}
\textbf{\te{లోతు}} — \emph{lōtu} — \textquotedblleft{}deep\textquotedblright{}.

\begin{grammarlens}[title={What you've built}]
One more word for this chapter.
\end{grammarlens}

\subsection*{Your turn}
\begin{itemize}
\raggedright
  \item \emph{Say it:} \emph{lōtu}
  \item \emph{Say it:} \emph{lōtu}, once more
  \item \emph{Say it:} say \emph{lōtu}
  \item \emph{Recall:} say the Telugu for fast; speed, then the Telugu for beautiful, then say \emph{lōtu} again
\end{itemize}

\begin{tcolorbox}[breakable,colback=teal!4,colframe=teal!35!black,title={Before you move on}]
What does \te{లోతు} mean? (\textquotedblleft{}Deep\textquotedblright{}.) Say it once more.
\end{tcolorbox}
\section[\texorpdfstring{\te{వెడల్పు} (veḍalpu)}{వెడల్పు (veḍalpu)}]{\te{వెడల్పు} (veḍalpu) — wide}

\label{lesson:TE-C138-vedalpu}



\begin{quote}
Before the new one: say the Telugu for beautiful, then the Telugu for deep.
\end{quote}

\subsection*{You'll want to know: \te{వెడల్పు}}
\textbf{\te{వెడల్పు}} — \emph{veḍalpu} — \textquotedblleft{}wide\textquotedblright{}.

\begin{grammarlens}[title={What you've built}]
One more word for this chapter.
\end{grammarlens}

\subsection*{Your turn}
\begin{itemize}
\raggedright
  \item \emph{Say it:} \emph{veḍalpu}
  \item \emph{Say it:} \emph{veḍalpu}, once more
  \item \emph{Say it:} say \emph{veḍalpu}
  \item \emph{Recall:} say the Telugu for beautiful, then the Telugu for deep, then say \emph{veḍalpu} again
\end{itemize}

\begin{tcolorbox}[breakable,colback=teal!4,colframe=teal!35!black,title={Before you move on}]
What does \te{వెడల్పు} mean? (\textquotedblleft{}Wide\textquotedblright{}.) Say it once more.
\end{tcolorbox}
\section[\texorpdfstring{\te{ఇరుకు} (iruku)}{ఇరుకు (iruku)}]{\te{ఇరుకు} (iruku) — narrow}

\label{lesson:TE-C138-iruku}



\begin{quote}
Before the new one: say the Telugu for deep, then the Telugu for wide.
\end{quote}

\subsection*{You'll want to know: \te{ఇరుకు}}
\textbf{\te{ఇరుకు}} — \emph{iruku} — \textquotedblleft{}narrow\textquotedblright{}.

\begin{grammarlens}[title={What you've built}]
One more word for this chapter.
\end{grammarlens}

\subsection*{Your turn}
\begin{itemize}
\raggedright
  \item \emph{Say it:} \emph{iruku}
  \item \emph{Say it:} \emph{iruku}, once more
  \item \emph{Say it:} say \emph{iruku}
  \item \emph{Recall:} say the Telugu for deep, then the Telugu for wide, then say \emph{iruku} again
\end{itemize}

\begin{tcolorbox}[breakable,colback=teal!4,colframe=teal!35!black,title={Before you move on}]
What does \te{ఇరుకు} mean? (\textquotedblleft{}Narrow\textquotedblright{}.) Say it once more.
\end{tcolorbox}
